% DMIR x microNAS — working draft
% Target: IEEE Signal Processing Letters (IEEEtran journal mode).
% Rules for this file:
%   * every number must come from logs/experiments.jsonl or a cited source;
%   * placeholders are \todo{...} — none may survive to submission;
%   * check prose against paper/STYLE.md before committing.
\documentclass[journal]{IEEEtran}

\usepackage{amsmath,amssymb}
\usepackage{booktabs}
\usepackage{graphicx}
\usepackage{xcolor}

\newcommand{\todo}[1]{\textcolor{red}{[TODO: #1]}}

\begin{document}

\title{Hardware-Constrained Architecture Search for Lane-Change Intention
Prediction from Vehicle Time Series}
% Alternative titles, pick later:
% - Microcontroller-Ready Lane-Change Intention Prediction via Constrained NAS
% - Searching Small: NAS under MCU Budgets for Driving Intention Time Series

\author{Sepehr~Mohammady,
        \todo{co-authors, order},
        \todo{supervisor}%
\thanks{The authors are with the ELIOS Lab, Department of Electrical,
Electronic and Telecommunications Engineering and Naval Architecture (DITEN),
University of Genoa, Genoa, Italy.}%
\thanks{This work was supported by the SYNERGIES project
\todo{grant number and official acknowledgment text}.}}

\maketitle

\begin{abstract}
Predicting a driver's lane-change intention early enough to act on it is a
time-series classification problem with hard deployment limits: the model has
to run on automotive-grade microcontrollers within tight memory and latency
budgets. We apply constrained neural architecture search to the Lane Change
Intention Recognition driving dataset \cite{lcdataset} (windows of 50 samples
over 31 signals) and obtain compact one-dimen-
sional convolutional networks for three tasks: three-class intention
classification and time-to-lane-change regression for left and right
maneuvers. \todo{results: accuracy, RMSE, model size, latency on STM32}
The searched models improve on the published baseline (Transformer,
0.51~s time-to-lane-change RMSE, deployed in FP32 \cite{lcintention2026})
while using \todo{X}$\times$ less memory, and we report measured, not
estimated, footprint and latency on the same STM32 targets after int8
quantization.
\end{abstract}

\begin{IEEEkeywords}
Neural architecture search, lane-change intention, time-series classification,
TinyML, embedded inference, STM32.
\end{IEEEkeywords}

\section{Introduction}
\todo{Motivation: intention prediction as advance warning for ADAS; why
on-device inference (latency, privacy, cost, no connectivity assumptions);
the gap: published DMIR models target accuracy, not deployment budgets.}

Contributions:
\begin{enumerate}
  \item a constrained search space and objective tuned to one-dimensional
        driving signals under microcontroller budgets;
  \item models that beat the published DMIR baseline on all three tasks at a
        fraction of its footprint \todo{verify once results exist};
  \item measured deployment numbers (flash, RAM, latency) on
        \todo{STM32 part}, with the full pipeline released as open source.
\end{enumerate}

\section{Related Work}
\subsection{Lane-change intention prediction}
\todo{The LC dataset and framework baseline \cite{lcintention2026} (TTLC
regression, Transformer RMSE 0.51~s, FP32 deployment on STM32H7B3/F401);
earlier HighD-based TTLC work (RMSE 0.63--0.7); MicroNAS for time series and
TinyTNAS as closest NAS relatives.}

\subsection{Constrained architecture search for microcontrollers}
\todo{muNAS \cite{liberis2021munas}, MCUNet/TinyNAS, MicroNets; what transfers
to 1D time series and what does not.}

\section{Data}
We use the Lane Change Intention Recognition dataset \cite{lcdataset}
(50 human drivers in the CARLA simulator, 10~Hz, over 3{,}400 annotated lane
changes; MIT license), prepared as fixed windows of $T{=}50$ samples
(5~s) over $F{=}31$ channels \todo{confirm channel order and whether channel
31 is fileTime — drop it if so; official feature count is 30}. Classification uses balanced splits of 94{,}620 / 16{,}332 /
19{,}722 windows (train/val/test) over three classes: no intention, right
lane change, left lane change. The regression tasks map a window to the time
until the lane change, $t \in [0, 4]$~s discretized at $0.1$~s, with
39{,}313 / 7{,}742 / 8{,}097 (LCL) and 33{,}201 / 5{,}828 / 7{,}073 (LCR)
windows.

The splits are user-independent: we verified by matching raw per-driver
session logs against the prepared windows (correlation over 50-sample
windows) that the validation and test drivers follow the protocol of
\cite{lcintention2026} with no window leakage across splits.
The test splits contain isolated spikes (up to $5\times10^{6}$) on the
right/left \texttt{curvatureDx} channel pair — a numerical-derivative
artefact absent from training; we clip validation and test features to the
per-channel training range and leave training data untouched.

\section{Method}
\subsection{Search space}
\todo{1D cells: depthwise-separable and standard convolutions, kernel sizes,
widths, depth, pooling; operator set restricted to what the ST Edge AI
compiler maps efficiently.}

\subsection{Constraints and search algorithm}
\todo{peak RAM, flash after int8 quantization, MACs; aging evolution with
constraint handling; search cost budget.}

\section{Experiments}
\subsection{Setup}
\todo{training schedule, hardware, seeds/repeats, hyperparameters; all runs
logged to the released experiments.jsonl.}

\subsection{Results}
Table~\ref{tab:results} reports test-set metrics. All searched-model numbers
are our own test-set evaluation of the saved models, from searches run to their
full 150-round budget. \todo{add peak-RAM, MACs, and measured on-device latency
columns from ST Edge AI; consider re-running with an all-models save policy for
a guaranteed-optimal front (see notes).}

\begin{table}[t]
\caption{Test-set comparison. The internal reference models (unpublished) were
evaluated on the same test set; classification accuracy is our own evaluation,
the Transformer RMSE values are as reported. Size is the parameter count.}
\label{tab:results}
\centering
\begin{tabular}{lcccc}
\toprule
Model & Params & Acc.\ (\%) & LCR & LCL \\
\midrule
Reference CNN (cls) & 441\,k & 91.5 & --- & --- \\
Reference Transf.\ (reg, RMSE) & 333/49\,k & --- & 0.42 & 0.44 \\
Our DSCNN baseline & 10\,k & 91.5 & 0.32$^\dagger$ & 0.33$^\dagger$ \\
\midrule
Ours searched (cls / reg MAE) & 84\,k & \textbf{92.1} & 0.29$^\dagger$ & 0.33$^\dagger$ \\
Ours compact & 8\,k & 91.3 & --- & --- \\
Ours, no turn-signal & 21\,k & 91.1 & --- & --- \\
\bottomrule
\end{tabular}
\\[2pt]
\footnotesize $^\dagger$MAE (our search optimizes MAE); reference regression is
reported as RMSE. Our LCR RMSE is 0.45 and LCL 0.50.
\end{table}

For classification the searched model is both smaller and more accurate than the
441\,k reference CNN (92.1\% at 84\,k parameters; a 8\,k model still reaches
91.3\%). Regression beats the published single-TTLC SOTA (RMSE 0.51) at 2--3$\times$
fewer parameters (our RMSE 0.45 LCR, 0.47 LCL). We do not claim to beat the
internal-reference Transformers on regression RMSE (0.42/0.44), which remain
competitive small models; an RMSE-aware search narrowed but did not close that
gap. Removing the turn-signal channels drops classification
accuracy by about one point (92.1\% to 91.1\%)---evidence that the model
anticipates from vehicle dynamics and surrounding traffic rather than reading a
declared signal.
\todo{RMSE columns and the internal-reference RMSE (0.42/0.44) are not yet
matched because the search optimizes MAE; decide RMSE-aware objective vs
MAE-primary framing.}

\subsection{Deployment on STM32}
Target platform: STM32H7B3I-DK discovery kit (Cortex-M7 at 280~MHz, 1.4~MB
SRAM, 2~MB flash) — the same MCU family used by the baseline framework
\cite{lcintention2026}, which deployed FP32 models only. We therefore compare in
FP32, like for like: our classifier is $5.1\times$ smaller than the reference
CNN (338~KB vs 1729~KB of weights) at higher accuracy (92.1\% vs 91.7\%), and a
$46\times$ smaller 37~KB model still reaches 91.3\%.

On quantization, the ST toolchain accepts only float32/int8/uint8 tensors. Full
int8 shrinks the classifier to 110~KB but costs accuracy on these
wide-dynamic-range inputs ($92.1\% \rightarrow 86.9\%$; LCR MAE
$0.287 \rightarrow 0.449$). An int16$\times$8 scheme (int8 weights, int16
activations) preserves accuracy almost exactly offline
($92.1\% \rightarrow 92.2\%$; LCR MAE $0.287 \rightarrow 0.286$), but the ST
toolchain does not support it---its documentation restricts weight and
activation tensors to float32, int8 and uint8---so we report it as an offline
bound rather than a deployed configuration.
\todo{cite ST Edge AI Core 4.0 documentation (quantization / supported ops,
retrieved 2026-07-14); confirm against the pending int16x8 benchmark.}
\todo{insert ST Edge AI measured latency / flash / peak SRAM per board
(STM32H7B3I-DK and NUCLEO-F401RE); see docs/research/deployment.md.}

\section{Conclusion}
\todo{two or three sentences; no new claims.}

\bibliographystyle{IEEEtran}
\begin{thebibliography}{9}
\bibitem{lcintention2026} L.~Forneris \emph{et al.}, ``A deployment-oriented
simulation framework for deep learning-based lane change prediction,''
\emph{IEEE Signal Process.\ Lett.}, vol.~33, pp.~136--140, 2026,
doi: 10.1109/LSP.2025.3638676.
\bibitem{liberis2021munas} E.~Liberis, {\L}.~Dudziak, and N.~D.~Lane,
``$\mu$NAS: Constrained neural architecture search for microcontrollers,'' in
\emph{Proc.\ 1st Workshop Mach.\ Learn.\ Syst.\ (EuroMLSys)}, 2021, pp. 70--79.
\bibitem{lcdataset} L.~Forneris \emph{et al.}, ``Lane change intention
recognition dataset,'' Zenodo, 2025, doi: 10.5281/zenodo.16686054.
\bibitem{forneris2024dmir} L.~Forneris \emph{et al.}, ``Setting up a
lightweight simulation environment for automated driving dataset
collection,'' in \emph{Applications in Electronics Pervading Industry,
Environment and Society (ApplePies)}, Springer LNEE, 2024, pp. 156--163,
doi: 10.1007/978-3-031-84100-2\_19.
\end{thebibliography}

\end{document}
